\documentclass{article}
\usepackage{amsmath,amssymb,amsthm}
\IfFileExists{hyperref.sty}{%
  \IfFileExists{color.sty}{\usepackage{color}\definecolor{paperlinkblue}{rgb}{0.10,0.25,0.42}%
    \usepackage[colorlinks=true,linkcolor=paperlinkblue]{hyperref}}{\usepackage[hidelinks]{hyperref}}%
}{}
\ifdefined\hyperlink
  \newcommand{\paperanchor}[1]{\hypertarget{#1}{}}%
  \newcommand{\paperlink}[2]{\hyperlink{#1}{\begin{NoHyper}#2\end{NoHyper}}}%
\else
  \ifdefined\pdfextension
    \ifnum\outputmode>0
      \def\paperdest{\pdfextension dest}\def\paperstart{\pdfextension startlink}%
      \def\paperend{\pdfextension endlink}\def\papercolor{\pdfextension colorstack}%
    \fi
  \else\ifdefined\pdfstartlink
    \ifnum\pdfoutput>0
      \let\paperdest\pdfdest\let\paperstart\pdfstartlink\let\paperend\pdfendlink\let\papercolor\pdfcolorstack
    \fi
  \fi\fi
  \ifdefined\paperdest
    \newcommand{\paperanchor}[1]{\paperdest name{#1} xyz}%
    \newcommand{\paperlink}[2]{\leavevmode\paperstart attr{/Border [0 0 0]} goto name{#1}%
      \papercolor 0 push {0.10 0.25 0.42 rg 0.10 0.25 0.42 RG}#2\papercolor 0 pop\paperend}%
  \else\ifdefined\XeTeXversion
    \newcommand{\paperanchor}[1]{\special{pdf:dest (#1) [@thispage /XYZ null @ypos null]}}%
    \newcommand{\paperlink}[2]{\leavevmode\special{pdf:bann << /Subtype /Link /Border [0 0 0] /A << /S /GoTo /D (#1) >> >>}%
      \special{color push rgb 0.10 0.25 0.42}#2\special{color pop}\special{pdf:eann}}%
  \else
    \PackageError{lean-to-latex}{No PDF hyperlink backend available}{Use XeLaTeX, pdfLaTeX, LuaLaTeX, or install hyperref.}%
  \fi\fi
\fi
\newcommand{\paperlineanchor}[1]{\raisebox{10pt}[0pt][0pt]{\paperanchor{#1}}}
\DeclareRobustCommand{\paperresultname}[1]{\paperlineanchor{theorem:\arabic{theorem}}#1}
\newtheorem{theorem}{\paperresultname{Theorem}}
\newcounter{paperpremise}
\begin{document}
\section*{Proofs in elementary number theory}

\subsection*{Prerequisites}
\begin{description}
\item[\paperlineanchor{premise:4e61742e74776f5f6d756c}\paperlink{premise:4e61742e74776f5f6d756c}{P1.}] \refstepcounter{paperpremise}\label{premise:4e61742e74776f5f6d756c}\textbf{Doubling by addition.} Let \(n\) be a natural number. Then \(2 \cdot n = n + n\).
\item[\paperlineanchor{premise:4e61742e6d756c5f7a65726f}\paperlink{premise:4e61742e6d756c5f7a65726f}{P2.}] \refstepcounter{paperpremise}\label{premise:4e61742e6d756c5f7a65726f}\textbf{Multiplication by zero.} Let \(n\) be a natural number. Then \(n \cdot 0 = 0\).
\item[\paperlineanchor{premise:4e61742e6476645f7a65726f}\paperlink{premise:4e61742e6476645f7a65726f}{P3.}] \refstepcounter{paperpremise}\label{premise:4e61742e6476645f7a65726f}\textbf{Divisibility of zero.} Let \(a\) be a natural number. Then \(a \mid 0\).
\item[\paperlineanchor{premise:4e61742e6476645f616464}\paperlink{premise:4e61742e6476645f616464}{P4.}] \refstepcounter{paperpremise}\label{premise:4e61742e6476645f616464}\textbf{Divisibility of a sum.} Let \(a\), \(b\) and \(c\) be natural numbers. If \(a \mid b\) and \(a \mid c\), then \(a \mid b + c\).
\item[\paperlineanchor{premise:4e61742e6d756c5f73756363}\paperlink{premise:4e61742e6d756c5f73756363}{P5.}] \refstepcounter{paperpremise}\label{premise:4e61742e6d756c5f73756363}\textbf{Multiplication by a successor.} Let \(n\) and \(m\) be natural numbers. Then \(n \cdot (m + 1) = n \cdot m + n\).
\item[\paperlineanchor{premise:4e61742e6d756c5f6173736f63}\paperlink{premise:4e61742e6d756c5f6173736f63}{P6.}] \refstepcounter{paperpremise}\label{premise:4e61742e6d756c5f6173736f63}\textbf{Multiplication associativity.} Let \(n\), \(m\) and \(k\) be natural numbers. Then \(n \cdot m \cdot k = n \cdot (m \cdot k)\).
\end{description}

\begin{theorem}\label{theorem:1}
Let \(n\) be a natural number. Then \(2 \mid n + n\).
\end{theorem}

\begin{proof}
Let \(n\) be a natural number.
By \paperlink{premise:4e61742e74776f5f6d756c}{Prerequisite~P\ref{premise:4e61742e74776f5f6d756c}}, we have \(2 \cdot n = n + n\)~(1).
Reversing the equality in (1) gives
\begin{equation}\tag{2}
n + n = 2 \cdot n.
\end{equation}
With witness \(n\), (2) shows that \(2 \mid n + n\).
\end{proof}

\begin{theorem}\label{theorem:2}
Let \(n\) be a natural number. If \(2 \mid n\), then \(\exists k \in \mathbb{N}\; (n = k + k)\).
\end{theorem}

\begin{proof}
Let \(n\) be a natural number.
Assume \(2 \mid n\)~(1).
By (1), choose a natural number \(k\) such that \(n = 2 \cdot k\)~(2).
By \paperlink{premise:4e61742e74776f5f6d756c}{Prerequisite~P\ref{premise:4e61742e74776f5f6d756c}}, we have \(2 \cdot k = k + k\)~(3).
Combining (2) and (3) gives \(n = k + k\)~(4).
With witness \(k\), (4) shows that
\[
\exists k_{1} \in \mathbb{N}\; (n = k_{1} + k_{1}).
\]
\end{proof}

\begin{theorem}\label{theorem:3}
Let \(a\) and \(n\) be natural numbers. Then \(2 \mid (a + a) \cdot n\).
\end{theorem}

\begin{proof}
Let \(a\) and \(n\) be natural numbers.
Record the instance of \paperlink{theorem:1}{Theorem~\ref{theorem:1}} with \(n = a\) as (1).
We proceed by induction on \(n\).
\begin{description}
\item[Base case.] We must show \(2 \mid (a + a) \cdot 0\).
By \paperlink{premise:4e61742e6d756c5f7a65726f}{Prerequisite~P\ref{premise:4e61742e6d756c5f7a65726f}}, we have \((a + a) \cdot 0 = 0\)~(2).
By \paperlink{premise:4e61742e6476645f7a65726f}{Prerequisite~P\ref{premise:4e61742e6476645f7a65726f}}, we have \(2 \mid 0\)~(3).
Rewriting (3) using (2) gives
\[
2 \mid (a + a) \cdot 0.
\]
\item[Inductive step.] Let \(m\) be a natural number.
As the induction hypothesis, assume
\begin{equation}\tag{4}
2 \mid (a + a) \cdot m.
\end{equation}
We must show \(2 \mid (a + a) \cdot (m + 1)\).
By \paperlink{premise:4e61742e6476645f616464}{Prerequisite~P\ref{premise:4e61742e6476645f616464}}, using (4) and (1), we have \(2 \mid (a + a) \cdot m + (a + a)\)~(5).
By \paperlink{premise:4e61742e6d756c5f73756363}{Prerequisite~P\ref{premise:4e61742e6d756c5f73756363}}, we have \((a + a) \cdot (m + 1) = (a + a) \cdot m + (a + a)\)~(6).
Rewriting (5) using (6) gives
\[
2 \mid (a + a) \cdot (m + 1).
\]
\end{description}
This completes the induction.
\end{proof}

\begin{theorem}\label{theorem:4}
Let \(a\), \(b\) and \(n\) be natural numbers. Then \(2 \mid (a + a) \cdot n + (b + b)\).
\end{theorem}

\begin{proof}
Let \(a\), \(b\) and \(n\) be natural numbers.
Record the instance of \paperlink{theorem:3}{Theorem~\ref{theorem:3}} with \(a = a\) and \(n = n\) as (1).
Record the instance of \paperlink{theorem:1}{Theorem~\ref{theorem:1}} with \(n = b\) as (2).
The claim follows from \paperlink{premise:4e61742e6476645f616464}{Prerequisite~P\ref{premise:4e61742e6476645f616464}}, using (1) and (2).
\end{proof}

\begin{theorem}\label{theorem:5}
Let \(n\) and \(m\) be natural numbers. If \(2 \mid n\), then \(2 \mid n \cdot m\).
\end{theorem}

\begin{proof}
Let \(n\) and \(m\) be natural numbers.
Assume \(2 \mid n\)~(1).
By (1), choose a natural number \(k\) such that \(n = 2 \cdot k\)~(2).
Applying \(k_{1} \mapsto k_{1} \cdot m\) to both sides of (2) gives
\begin{equation}\tag{3}
n \cdot m = 2 \cdot k \cdot m.
\end{equation}
By \paperlink{premise:4e61742e6d756c5f6173736f63}{Prerequisite~P\ref{premise:4e61742e6d756c5f6173736f63}}, we have \(2 \cdot k \cdot m = 2 \cdot (k \cdot m)\)~(4).
Combining (3) and (4) gives
\begin{equation}\tag{5}
n \cdot m = 2 \cdot (k \cdot m).
\end{equation}
With witness \(k \cdot m\), (5) shows that \(2 \mid n \cdot m\).
\end{proof}

\end{document}
